\section{A logarithmic-depth ordinary-tanh normal form}
\label{app:normal}

We supply the geometric reduction and all remaining normalization
steps. They are needed because the estimator acts on tanh itself,
whereas a generic finite-precision simulation could replace tanh by
lookup tables without preserving the hypothesis of the estimator.

\subsection{Relevant coordinates and a two-ball cover}

\begin{lemma}[Relevant-coordinate reduction]
\label{lem:relevant}
A globally unit-margin tanh representation has $P\ge1$, and its
Boolean function has at most $P^2$ relevant input coordinates.
Fixing all the other coordinates can be done without increasing its
product or reducing its global margin on the remaining cube.
\end{lemma}
\begin{proof}
Since tanh is zero-preserving and one-Lipschitz,
\[
 1\le\E_x z(x)^2\le
 \left(\prod_{s=2}^D\norm{W^s}_F^2\right)
 \E_x\norm{W^1(1,x)}_2^2=P^2,
\]
where the expectation is uniform and
$\E[(1,x)(1,x)^T]=I$. Put $T=\prod_{s=2}^D\norm{W^s}_F$.
For every relevant coordinate $i$, a bichromatic cube edge in that
direction gives $2\le2T\norm{W^1e_i}_2$. Summing squares yields
at most $P^2$ relevant coordinates. For the irrelevant columns,
choose independent uniform signs $a_i$, replace the first constant
column $w_0$ by $w_0+\sum_i a_iw_i$, and delete those columns.
The expected new squared Frobenius norm equals the old squared norm.
Some choice does not increase it. Every restriction has the same
Boolean labels and retains the original margin.
\end{proof}

\begin{lemma}[Two balls]
\label{lem:twoball}
For every real matrix $W$, every Euclidean ball $B(a,r)$, and every
coordinatewise one-Lipschitz map $\Phi$, the set
$\Phi(W B(a,r))$ is covered by two balls of radius
$r\norm W_F/\sqrt2$.
\end{lemma}
\begin{proof}
Translation and scaling reduce to $a=0,r=1,\norm W_F=1$;
the zero case is immediate. Let $\alpha$ be the leading singular
value. If $\alpha^2\le1/2$, one ball suffices. Otherwise put
$b^2=1-\alpha^2$ and $c=(2\alpha^2-1)/(2\alpha)$.
The remaining singular part has operator norm at most $b$.
With $u$ a leading left singular vector, a point whose leading
input coordinate is $t$ has squared distance to the matching
center $cu$ or $-cu$ at most
\[
 (\alpha|t|-c)^2+b^2(1-t^2).
\]
This is convex in $|t|\in[0,1]$ and both endpoint values equal
$1/(4\alpha^2)\le1/2$. Apply $\Phi$ to the two centers;
its Euclidean Lipschitz constant is at most one.
\end{proof}

For prefix products $S_j=\prod_{s\le j}\norm{W^s}_F$ and
inputs of norm at most $R$, iteration gives at most $2^j$ balls of
radius $RS_j2^{-j/2}$. The remaining suffix is Lipschitz with
constant at most $P/S_j$. Thus normalized prefix states of opposite
labels are at least $2/P$ apart. If
\[
 j=\left\lceil2\log_2(64PR)\right\rceil,
\]
selecting one actual state from each nonempty ball gives polynomially
many centers $c_a=h^j(x_a)/S_j$, each with its label $f_a$.
Every normalized state has a same-label center within $1/(8P)$,
and every opposite-label center is at least $2/P$ away. All states
and centers have norm at most $R$. The center need not initially lie
in the image: selecting an actual image doubles the cover radius,
which the displayed constants allow. Every covering ball is
monochromatic by the suffix Lipschitz estimate.

\subsection{Physical weights and finite precision at short depth}

\begin{lemma}[Bounding physical weights]
\label{lem:physical}
Suppose a globally unit-margin tanh representation already has
depth $D$, hidden widths at most $C$, and product at most $P$.
It has a representation of the same Boolean function with no more
weighted layers, product at most $2P$, margin greater than $3/2$,
and individual weights bounded by a universal polynomial in
$d,D,C,P$. It can then be rounded to dyadic weights with
$O(D\log B)$ bits and product at most $4P$, where $B$ is a
polynomial bounding its dimensions, Frobenius norms, and row sums.
The rounded representation retains unit margin.
\end{lemma}
\begin{proof}
Write $M_j=\norm{W^j}_F$, $S_j=\prod_{s\le j}M_s$, and
$R=\sqrt{d+1}$. Margin and bounded hidden activations give
$S_j\le P\sqrt C$ at every hidden layer. Set
\[
 \tau^2=\frac1{16PD R^3}.
\]
Replace tanh by the identity at each hidden layer whose original
prefix product is less than $\tau$. For any intermediate hybrid
network, all prefix and suffix Lipschitz bounds remain valid.
Using $|u-\tanh u|\le|u|^3/3$ and
$\norm{v^{\odot3}}_2\le\norm v_2^3$, the telescoping output
error is at most
\[
 \sum_{j:S_j<\tau}\frac{P}{S_j}\frac{(S_jR)^3}{3}
 \le PD\tau^2R^3/3\le1/48.
\]
Collapse consecutive linear blocks into the following nonlinear
layer or final readout, then double the readout. If retained hidden
indices are $j_1<\cdots<j_H$, the collapsed matrices have assigned
Frobenius bounds
$S_{j_1}$, $S_{j_s}/S_{j_{s-1}}$, and $2P/S_{j_H}$.
Since each retained $S_{j_s}$ lies between $\tau$ and $P\sqrt C$,
all weights are bounded by $\max(1,2P\sqrt C/\tau)$, a polynomial.
The product is at most $2P$ and the margin exceeds $3/2$.
If no hidden activation remains, the affine matrix has norm at
most $2P$ and the same conclusions hold.

For precision, enlarge a polynomial $B\ge2$ to bound all dimensions,
matrix Frobenius norms, and absolute row sums. At remaining depth
$D'\le D$, every nonzero matrix norm satisfies
\[
 M_j\ge1/(\sqrt{d+1}B^{D'-1}),
\]
because $1\le|z(x)|\le\sqrt{d+1}\prod_jM_j$.
Round entries to mesh $2^{-b}$ with
$b=O(D'\log B+\log(d+D'+2))$, with a sufficiently large
universal constant. Matrix perturbations are at most $M_j/(4D')$,
so the product increases by less than two. The layerwise error
recursion using the row-sum bound $B$ and Lipschitzness of tanh
makes the output error less than $1/2$ at the same precision.
The pre-rounding margin in fact exceeds $2(1-1/48)>3/2$, so the
remaining margin is greater than one. If needed, absorb the additive
logarithm into $O(D\log B)$ by enlarging $B$.
\end{proof}

\subsection{Replacing the long suffix by a tanh classifier}

\begin{proof}[Proof of Theorem~\ref{thm:normal}]
Apply Lemma~\ref{lem:relevant}, retain $k\le P^2$ inputs, and
put $R=\sqrt{k+1}$ and
$j=\lceil2\log_2(256PR)\rceil$.
If $D=1$, the affine coefficients have norm at most $P$; double
the readout and round them directly to sufficiently fine dyadic
precision $O(\log(P+2))$, retaining unit margin, and finish.
Hence assume $D\ge2$.
Use the prefix version of the compression theorem in
Appendix~\ref{sec:pf-width-compression} on the first
$\min(j,D-1)$ hidden layers with error $1/4$, then double the
readout. The global margin is at least $3/2$, the full product is
at most $P_0=4P$, and all retained widths are bounded by a
universal polynomial $C=C(P)$: the only inputs to that width
bound are $P,R,j,\log(2|X_k|)$ and a fixed error tolerance.
If all hidden layers have been retained, apply
Lemma~\ref{lem:physical} and finish. Otherwise the first $j$
hidden layers have the cover and normalized centers just proved,
with $P_0$ in place of $P$.

First supply a constant to a new suffix while preserving the product
bound. For retained prefix products $S_s$, global margin gives
$S_s\le P_0\sqrt C$. Choose
\[
 \epsilon_0=1/\max(1,2\max_{s\le j}S_s).
\]
It has inverse polynomial in $P$ and satisfies
$0<\epsilon_0S_s\le1/2$. Add one constant rail whose activation
in layer $s$ is exactly $\epsilon_0S_s$. Its first weight is
$\operatorname{atanh}(\epsilon_0S_1)$ on the augmented constant;
its later weight is
$\operatorname{atanh}(\epsilon_0S_s)/(\epsilon_0S_{s-1})$.
For $0<u\le1/2$,
$(\operatorname{atanh}u/u)^2\le1+u^2\le5/4$.
Thus the added squared norm is at most $(5/4)M_s^2$ in each
layer, and the augmented prefix product is at most $(3/2)^jS_j$.

At the interface the features are $(h^j,\epsilon_0S_j)$.
A suffix written for $(u,1)=(h^j/S_j,1)$ is attached by multiplying
its first matrix by
$\operatorname{diag}(S_j^{-1}I,(\epsilon_0S_j)^{-1})$.
This increases its first Frobenius norm by at most
$(\epsilon_0S_j)^{-1}$. The complete product is therefore at
most $(3/2)^j\epsilon_0^{-1}$ times the suffix's ordinary product.
Both extra factors are polynomial in $P$.

We now give the suffix. Its input $u$ and each of the $q=\poly(P)$
centers $c_a$ have norm at most $R$. A correct-label center is
within $1/(8P_0)$, while each opposite-label center is at distance
at least $2/P_0$. In one hidden layer provide the features
\[
 \tanh(\beta u_i),\qquad
 \tanh(1+\alpha u_i),\qquad\tanh(1-\alpha u_i),
\]
and an exact constant $1/2$. Linear combinations of these features
approximate respectively $u_i$ and $u_i^2$ using
\[
 \frac{\tanh(\beta u_i)}\beta,\qquad
 \frac{\tanh(1+\alpha u_i)+\tanh(1-\alpha u_i)-2\tanh1}
      {\alpha^2\tanh''1}.
\]
Their errors are $O(\beta^2R^3)$ and $O(\alpha^2R^4)$.
Choose $\alpha,\beta$ inverse-polynomially small in
$C,R,P_0$ so that every affine combination
\[
 D_a(u)=\sum_i\operatorname{approx}(u_i^2)
       -2\sum_i(c_a)_i\operatorname{approx}(u_i)+\norm{c_a}_2^2
\]
approximates $\norm{u-c_a}_2^2$ to error at most
$1/(256P_0^2)$. This is possible with polynomial weights because
$\norm{c_a}_1\le\sqrt C R$ and $\tanh''1\ne0$.
Every affine constant is charged to the constant rail.

In a second hidden layer use
\[
 t_a(u)=\tanh(B_0[\tau_0-D_a(u)]),\qquad
 \tau_0=1/(16P_0^2),\quad B_0=32P_0^2\log(16q),
\]
and propagate the constant $1/2$. Put $\chi_a=(1+t_a)/2$.
For a selected close center, $\tau_0-D_a\ge11/(256P_0^2)$;
for every opposite-label center it is less than $-3/P_0^2$.
The tanh tail estimate gives $\chi_a\ge1-1/(16q)$ for a close
center and $\chi_a\le1/(16q)$ for each opposite-label center.
All other $\chi_a$ lie in $[0,1]$, so
\[
 f(u)\sum_a f_a\chi_a(u)\ge1-1/(16q)-q/(16q)\ge7/8.
\]
The readout $2\sum_af_a\chi_a=\sum_af_at_a+\sum_af_a$
has margin at least $7/4$; implement its constant by the carried
$1/2$ rail. This is an ordinary-tanh network with two hidden
layers, polynomial width, polynomial weights, and polynomial product.

Attach it to the augmented prefix. The resulting network has
depth $j+3=O(\log(P+2))$ and polynomial width and product.
Some interface weights can be large if $S_j$ is tiny. Apply
Lemma~\ref{lem:physical} to this now-short network. It preserves
ordinary tanh and the Boolean function, makes all physical weights
polynomial in $P$, and rounds them with
$O(D'\log B)=O(\log^2(P+2))$ bits, since both the new depth
$D'$ and $\log B$ are $O(\log(P+2))$.
This proves all parts of the normal form.
\end{proof}
